\Qhead{4.3}{Режимы работы и шумы лазеров}{Режимы работы лазеров: непрерывный режим генерации, режим модуляции добротности резонатора, режим синхронизации мод. Сверхкороткие импульсы. Шумы лазеров, формула Таунса и предельная стабильность частоты.}
\Subfile{4.3.1}{sections/04_generation/q03/s01}
\Subfile{4.3.2}{sections/04_generation/q03/s02}

\recap{Q-модуляция высвобождает накопленную инверсию; синхронизация мод складывает широкую полосу с определёнными фазами. Фазовая диффузия задаёт фундаментальную ширину, технические шумы обычно увеличивают её.}
\board{Временная последовательность Q-затвора; сумма фазированных мод; экспоненциальная когерентность и лоренцева линия.}
\oral{Почему устойчивость мощности не фиксирует фазу свободного лазера?}
\bait{Можно оставить различие времени жизни поля и энергии: именно оно часто создаёт множитель два при сравнении формул ширины.}
\sources{Для коэффициента лазерной линии: Rana~\cite{rana}, PDF-слайд~37; расхождение вдвое с записью Kärtner (полный PDF-с.~327, печатная с.~319) пока не согласовано. Kärtner~\cite{kartnerfull}, главы~4, 5, 7 и 9 (режимы и шумы); университетские слайды по лазерам~\cite{optoelectronics}. В \eqref{eq:schawlow} все скорости и ширины определены явно; универсальное применение коэффициента к любому лазеру не предполагается.}
